% !TeX root = ../paper.tex

\section{Apparatus}
\label{sec:apparatus}

The player starts inside a dimly lit wooden cabin. A lock with four coin-shaped slots bars the front door; the way out is to find four golden coins hidden around the room and insert them into the lock. Dark ambient music and the cramped, cluttered interior set the tone before anything threatens the player.

Each coin hides inside one of four challenges, described in Section~\ref{sec:apparatus-challenges}, so the search forces the player to interact with every station in the room. The game has no time limit and no fail state; players can take as long as they want, and nothing punishes hesitation. The discomfort lies in what the player has to touch to get each coin. Once the fourth coin clicks into place, the lock opens and a win screen ends the session. We chose this flow for the study because it guarantees that each participant encounters all four haptic (or non-haptic) interactions in a single playthrough, while leaving the order of the challenges up to the player.

A game manager script drives this loop. The same scene runs in both study conditions; the conditions differ in whether the physical proxies lie in the tracking space at the positions of their virtual counterparts.

% !TeX root = ../../paper.tex

\subsection{The Four Challenges}
\label{sec:apparatus-challenges}

Each challenge pairs a virtual horror object with an everyday physical proxy chosen to match its expected texture. The proxies cost a few euros each; none of them is instrumented.

\paragraph{Bowl of eyes.} A rusted bowl on a wooden table is filled with realistic looking eyeballs. The player has to dig through them to find the coin at the bottom of the bowl. The physical proxy is a real metal bowl filled with rubber balls whose size matches the virtual eyes exactly; the virtual bowl is scaled at runtime to the measured dimensions of the physical one, so the rim is where the hand expects it.

\paragraph{Bucket of maggots.} A wooden bucket teems with writhing maggots, again with a coin buried underneath. The proxy is a bucket of gummy worms. Their damp, slightly sticky surface and the way a handful gives under pressure come close to what the visuals promise, at least closer than any dry material we tried.

\paragraph{Bloody pipe.} A pipe assembly bursts when the player pulls out the coin and sprays blood at them, driven by a particle system. The haptic counterpart is deliberately low-tech, in the tradition of Wizard-of-Oz setups: a researcher standing outside the play area sprays a fine mist of water from a bottle through a plastic tube aligned with the virtual pipe, timed to the burst. The spray lasts around five seconds and then stops together with the particle effect.

\paragraph{Cobwebs behind the closet.} The last coin hides behind a closet, in a mass of spider webs the player has to reach into. The proxy is synthetic Halloween spider webbing stretched behind an upright table. The material clings to skin and stretches when the hand pulls back---a sensation that is hard to fake with anything else and unpleasant even when one knows exactly what it is. The coin itself is hanging on a thin string, so the player feels the sensation of ripping it free when they attempt to remove it.

The set covers a deliberate range of contact types: immersing the whole hand (eyes, worms), a sudden stimulus arriving on the body (spray), and a clinging texture the player has to pull free from (webs). It also varies who controls the moment of contact---the player in three challenges, the experimenter in one.

% !TeX root = ../../paper.tex

\subsection{Implementation and Alignment}
\label{sec:apparatus-implementation}

The prototype runs standalone on a Meta Quest~3 and is built in Unity~6 with the Universal Render Pipeline, using the Meta XR SDK via OpenXR. The cabin interior, props, and creatures are assembled from asset store packs; custom scripts handle everything specific to the study.

Getting a proxy illusion to hold comes down to alignment, and we use two mechanisms depending on whether an object moves. Static furniture and the proxy stations are aligned by hand during setup: the experimenter places the physical bowl, bucket, pipe outlet, and webbing in the room, adjusts the corresponding virtual objects to match, and pins the result with the Quest's spatial anchors so the calibration survives headset restarts and short tracking losses. A dedicated authoring mode supports this step, and a world-lock component keeps each virtual object fixed to its anchored real-world pose at runtime.

The coins are the one thing the player carries around, so a static anchor is not enough. Physical coin props are tracked optically through the headset's passthrough cameras using AprilTag fiducial markers, with a One-Euro filter smoothing the raw poses. The virtual coin follows its physical twin, which lets a participant fish a coin out of the maggot bucket, hold onto it, and slot it into the door lock, with the virtual and physical coin staying in register throughout.

Scale mismatches are handled at the source: a sizing component scales each virtual object so that its rendered dimensions match the measured dimensions of its physical proxy in centimeters, either uniformly or per axis. This mattered more than expected---a virtual bowl a few centimeters wider than the real one is enough for the player's fingers to visibly clip through the rim on approach.

